\chapter{Outlook}\label{ch:outlook}

This book has followed WRF and WRF-Fire from their equations to their port onto GPUs. The port's guiding decision was
to accept no difference at all: a GPU build that gives the same bits as the CPU build, so that every difference is a
bug and every bug can be found. This last chapter looks at what that decision makes possible next.

\section{The new fire model}

WRF~4.8.0 brings NCAR's Community Fire Behavior Model (\cref{ch:cfbm}). Its algorithms are those of WRF-Fire, rewritten
as a modular, standalone code, so its GPU port can reuse the kernels of WRF-Fire's: level-set stencils, the spread
rate per node, fuel consumption, the flux sums. Its standalone driver allows the port to be developed and tested
without WRF, on a small grid, in seconds. The method stays the same: bitwise equality with the CPU build of the same
source, standalone first, then inside WRF~4.8.0. And since CFBM is under active development, the GPU changes should be
contributed to NCAR's repository rather than kept in a fork.

\section{Other GPUs}

The port targets NVIDIA's A100 and H100, by the project owner's choice of speed over portability (\cref{ch:gpu}).
Porting to AMD GPUs would mean rewriting the OpenACC directives as OpenMP offload, which AMD's compilers support,
and the few CUDA Fortran kernels in HIP. The verification would carry over unchanged: contraction off, the portable
math library, and the same tests. And since REPRO uses only IEEE basic operations and the portable functions, the
same bits on every vendor's GPU are within reach. Different hardware would then compute the same fire.

\section{The FAST build}

Users who run forecasts want speed more than bit equality with a reference. The FAST build mode (ADR-003) uses the
same source and the same kernels with the vendors' fast arithmetic: fused multiply-add, vendor math libraries, and,
behind macros, the parallel floating-point reductions that REPRO forbids. Its results will differ from REPRO's in
the last bits, and soon by more, since the fire amplifies differences (\cref{ch:repro}). So FAST is validated
statistically: an ensemble of REPRO runs with perturbed initial states defines the spread a correct model may show,
and a FAST run must fall inside it, in the manner of the ensemble consistency tests of \citet{baker2015} and
\citet{milroy2018}. The fire perimeter is compared against the ensemble's envelope. The speed difference between the
two modes is itself a result: the measured cost of reproducibility.

\section{More of WRF}

Stage~8 of the plan extends the port beyond the case's configuration, one option family per phase: restarts on the
GPU, two-way nesting, more than two domains, the adaptive time step, other advection and turbulence options,
Thompson, Morrison and TEMPO microphysics, MYNN, RRTMG shortwave, Noah-MP, cumulus schemes for coarse outer domains,
and the other fire options. Each phase brings its own small case, its own CPU reference and its own bitwise gate. The
order follows what wildfire studies use.

\section{Upstream}

A port maintained apart from NCAR's code must be carried forward with every release (\cref{ch:versions}). The lasting
solution is upstream: directives in NCAR's own source, under a macro that leaves the CPU build untouched, as the port
already arranges it. The verification tools (the bit-hash tracer, the call check, the reproducible math library)
would be useful there too, for any change to WRF, GPU or not: they turn ``the results look similar'' into ``the
results are identical'' or ``this field, at this step, in this routine, differs''.
